% 2-3-classical-ml.tex
%  Budget: 2pp.
%  Rules: every numeral in \lr{}; cite only from references.bib;
%  em-dash must be \lr{---} (bare --- is a tofu box in B Nazanin).

\section{یادگیری ماشین کلاسیک}

خانواده دوم، روش‌های یادگیری ماشین پیش از عصر یادگیری عمیق است: شبکه‌های
عصبی کم‌عمق، ماشین‌های بردار پشتیبان، روش‌های مبتنی بر درخت و روش‌های
ترکیبی. مدل \lr{LightGBM} در پژوهش حاضر به همین خانواده تعلق دارد.

فروزان و همکاران \cite{foruzan_comparative_2015} مطالعه‌ای مقایسه‌ای میان
روش‌های مختلف یادگیری ماشین انجام می‌دهند و مقدمه آنان نکته‌ای را بیان
می‌کند که در سراسر این ادبیات تکرار می‌شود: پیش‌بینی دقیق قیمت برق ذاتاً
دشوار است، اما همین دقت است که به تولیدکنندگان و مصرف‌کنندگان امکان
می‌دهد پیشنهاد راهبردی بدهند و ریسک خود را کاهش دهند.

کاستلی و همکاران \cite{castelli_forecasting_2020} رویکردی مبتنی بر
برنامه‌نویسی ژنتیک را برای بهبود دقت پیش‌بینی به کار می‌گیرند و بر
داده‌های تجربی یکی از بزرگ‌ترین بازارها تکیه می‌کنند. شاه و همکاران
\cite{shah_short-term_2021} از تکنیک‌های یادگیری ماشین گروهی\LTRfootnote{Ensemble}
برای پیش‌بینی کوتاه‌مدت قیمت استفاده می‌کنند و اهمیت آن را از منظر
مدیریت بهینه شبکه و برای بازیگران بازار \lr{---} شرکت‌های تولیدکننده،
مصرف‌کنندگان و بهره‌بردار سیستم \lr{---} برجسته می‌سازند. حضور همین ایده
ترکیب در ادبیات، یکی از انگیزه‌های پرسش دوم پژوهش حاضر است.

سان و همکاران \cite{sun_day-ahead_2022} راهبردی مبتنی بر یادگیری ماشین و
بهینه‌سازی برای پیش‌بینی قیمت روز-پیش ارائه می‌کنند و ضریب بار سیستم را
به‌عنوان متغیر ورودی وارد می‌کنند. الکواز و همکاران
\cite{alkawaz_day-ahead_2022} مدلی مبتنی بر رگرسیون ترکیبی پیشنهاد می‌دهند و
سه ویژگی چالش‌برانگیز قیمت را صریحاً فهرست می‌کنند: نوسان‌پذیری بالا، جهش
سریع، و فصلی‌بودن. یوساف و همکاران \cite{yousaf_novel_2021} روشی برای
سامانه‌های مدیریت انرژی ارائه می‌کنند که در آن متغیرهایی مانند قیمت خرید
ظرفیت، قیمت خرید توان، نرخ تعرفه و تقاضای بار به کار گرفته می‌شوند.

ناز و همکاران \cite{naz_short-term_2019} در چارچوب شبکه هوشمند، پیش‌بینی
همزمان بار و قیمت را با ماشین یادگیری فرین بهبودیافته بررسی می‌کنند.
البهلی و همکاران \cite{albahli_electricity_2020} کاربردی متفاوت را هدف
می‌گیرند: پیش‌بینی قیمت برق برای مراکز داده و رایانش ابری، که مصرف برق
بخش عمده هزینه عملیاتی آن‌هاست.

\subsection*{شواهد مربوط به شرایط غیرعادی بازار}

دو مطالعه در این خانواده، به‌طور خاص به پژوهش حاضر مربوط‌اند، زیرا به رفتار
مدل‌ها در شرایطی می‌پردازند که با دوره آموزش تفاوت دارد.

آریا و همکاران \cite{arya_machine_2021} پیش‌بینی قیمت برق را در دوره
همه‌گیری \lr{COVID-19} بررسی می‌کنند و یادآور می‌شوند که این دوره،
قیمت‌گذاری قابل اتکا را به مسئله‌ای پیش‌بینی‌ناپذیر بدل کرد و آماده‌سازی
برای پیش‌بینی آینده را دشوار ساخت. هالوژان و همکاران
\cite{haluzan_performance_2020} نیز عملکرد الگوریتم‌های جایگزین را بر
شبیه‌سازی بازارهای یونان و مجارستان و بر یک سری زمانی بلند، از ژانویه
\lr{2015} تا سپتامبر \lr{2018}، ارزیابی می‌کنند.

هر دو مطالعه به این نکته اشاره دارند که عملکرد یک مدل، تابعی از شرایط
بازاری است که در آن ارزیابی می‌شود. پژوهش حاضر همین ملاحظه را یک گام
جلوتر می‌برد: به‌جای گزارش عملکرد در یک دوره، مدل‌های تثبیت‌شده را بر
دوره‌ای کاملاً متفاوت می‌آزماید و سازوکار افت عملکرد را نیز اندازه‌گیری
می‌کند.
